% Tabla A2-T1. Generada automáticamente; no editar a mano.
% Requiere \usepackage{booktabs} (y graphicx si la tabla es ancha).
\begin{table}[htbp]
    \centering
    \caption{Análisis II: construcción de la muestra. Número de planetas y de estrellas después de cada corte sucesivo sobre el catálogo limpio.}
    \label{tab:A2-T1}
    \small
    \begin{tabular}{p{0.40\textwidth}rr}
    \toprule
    Corte & $N_{\rm pl}$ & $N_\star$ \\
    \midrule
    Masa y radio disponibles & 6\,199 & 4\,676 \\
    Sin límites superiores o inferiores & 6\,013 & 4\,598 \\
    Solo masa verdadera (pl\_bmassprov = Mass) & 2\,165 & 1\,793 \\
    Radio medido (planeta observado en tránsito) & 1\,504 & 1\,224 \\
    Con masa estelar & 1\,504 & 1\,224 \\
    Con metalicidad & 1\,488 & 1\,208 \\
    Con flujo incidente calculado & 1\,488 & 1\,208 \\
    Planetas pequeños (1 $<$ Rp $<$ 4) & 574 & 385 \\
    Masa no mayor que la de un planeta de hierro puro (Zeng et al. 2016) & 535 & 365 \\
    \bottomrule
    \end{tabular}
    \par\smallskip
    {\footnotesize\textit{Nota:} Composición final por régimen de radio: terrestre 31, super-tierra 142, tipo neptuno 362.}
\end{table}
